\begin{figure*}[t]
\centering
\begin{tikzpicture}
\draw[ge] (1.289,4.259) -- (2.017,3.605);
\draw[ge] (1.289,4.259) -- (3.524,3.932);
\draw[ge] (1.289,4.259) -- (2.329,4.750);
\draw[ge] (1.289,4.259) -- (1.289,2.295);
\draw[ge] (2.017,3.605) -- (4.200,2.459);
\draw[ge] (2.017,3.605) -- (3.524,3.932);
\draw[ge] (2.017,3.605) -- (2.329,4.750);
\draw[ge] (2.017,3.605) -- (1.289,2.295);
\draw[ge] (4.200,2.459) -- (3.524,3.932);
\draw[ge] (4.200,2.459) -- (2.641,2.295);
\draw[ge] (3.524,3.932) -- (2.329,4.750);
\draw[ge] (3.524,3.932) -- (2.641,2.295);
\draw[ge] (3.524,3.932) -- (3.524,0.659);
\draw[ge] (2.641,2.295) -- (3.524,0.659);
\draw[ge] (2.641,2.295) -- (1.289,2.295);
\draw[ge] (3.524,0.659) -- (1.289,2.295);
\draw[ge] (1.289,2.295) -- (0.666,1.641);
\draw[ge] (0.458,0.659) -- (0.666,1.641);
\draw[ge] (0.458,0.659) -- (1.082,0.250);
\draw[ge] (0.250,2.459) -- (0.666,1.641);
\node[ring] at (2.641,2.295) {};
\node[gn, fill=pc0, text=white] (P1n1) at (1.289,4.259) {1};
\node[gn, fill=pc0, text=white] (P1n2) at (2.017,3.605) {2};
\node[gn, fill=pc0, text=white] (P1n3) at (4.200,2.459) {3};
\node[gn, fill=pc0, text=white] (P1n4) at (3.524,3.932) {4};
\node[gn, fill=pc0, text=white] (P1n5) at (2.329,4.750) {5};
\node[gn, fill=pc1, text=white] (P1n6) at (2.641,2.295) {6};
\node[gn, fill=pc1, text=white] (P1n7) at (3.524,0.659) {7};
\node[gn, fill=pc1, text=white] (P1n8) at (1.289,2.295) {8};
\node[gn, fill=pc1, text=white] (P1n9) at (0.458,0.659) {9};
\node[gn, fill=pc1, text=white] (P1n10) at (0.250,2.459) {10};
\node[gn, fill=pc1, text=white] (P1n11) at (0.666,1.641) {11};
\node[gn, fill=pc1, text=white] (P1n12) at (1.082,0.250) {12};
\node[ptitle] at (2.225,5.000) {(b) Stage 2 with Jaccard (TPOC-LPA)};
\node[pstat] at (2.225,0.000) {2 comm., NMI 0.362};
\node[barrow, rotate=180, minimum height=9mm] at (5.100,2.500) {};
\node[note, anchor=south] at (5.100,2.780) {Jaccard};
\draw[ge] (7.605,4.259) -- (8.729,3.605);
\draw[ge] (7.605,4.259) -- (11.057,3.932);
\draw[ge] (7.605,4.259) -- (9.211,4.750);
\draw[ge] (7.605,4.259) -- (7.605,2.295);
\draw[ge] (8.729,3.605) -- (12.100,2.459);
\draw[ge] (8.729,3.605) -- (11.057,3.932);
\draw[ge] (8.729,3.605) -- (9.211,4.750);
\draw[ge] (8.729,3.605) -- (7.605,2.295);
\draw[ge] (12.100,2.459) -- (11.057,3.932);
\draw[ge, line width=1.5pt, draw=black!55] (12.100,2.459) -- (9.692,2.295);
\draw[ge] (11.057,3.932) -- (9.211,4.750);
\draw[ge, line width=1.5pt, draw=black!55] (11.057,3.932) -- (9.692,2.295);
\draw[ge] (11.057,3.932) -- (11.057,0.659);
\draw[ge, line width=1.5pt, draw=black!55] (9.692,2.295) -- (11.057,0.659);
\draw[ge, line width=1.5pt, draw=black!55] (9.692,2.295) -- (7.605,2.295);
\draw[ge] (11.057,0.659) -- (7.605,2.295);
\draw[ge] (7.605,2.295) -- (6.642,1.641);
\draw[ge] (6.321,0.659) -- (6.642,1.641);
\draw[ge] (6.321,0.659) -- (7.284,0.250);
\draw[ge] (6.000,2.459) -- (6.642,1.641);
\path (9.692,2.295) -- (12.100,2.459) node[elab, sloped, pos=0.45] {0.17\,|\,0.50};
\path (9.692,2.295) -- (11.057,3.932) node[elab, sloped, pos=0.60] {0.25\,|\,0.67};
\path (9.692,2.295) -- (11.057,0.659) node[elab, sloped, pos=0.50] {0.40\,|\,0.67};
\path (9.692,2.295) -- (7.605,2.295) node[elab, sloped, pos=0.50] {0.12\,|\,0.50};
\node[ring] at (9.692,2.295) {};
\node[gn, fill=pc0, text=white] (P2n1) at (7.605,4.259) {1};
\node[gn, fill=pc0, text=white] (P2n2) at (8.729,3.605) {2};
\node[gn, fill=pc0, text=white] (P2n3) at (12.100,2.459) {3};
\node[gn, fill=pc0, text=white] (P2n4) at (11.057,3.932) {4};
\node[gn, fill=pc0, text=white] (P2n5) at (9.211,4.750) {5};
\node[gn, fill=pc1, text=white] (P2n6) at (9.692,2.295) {6};
\node[gn, fill=pc1, text=white] (P2n7) at (11.057,0.659) {7};
\node[gn, fill=pc1, text=white] (P2n8) at (7.605,2.295) {8};
\node[gn, fill=pc3, text=black] (P2n9) at (6.321,0.659) {9};
\node[gn, fill=pc2, text=black] (P2n10) at (6.000,2.459) {10};
\node[gn, fill=pc1, text=white] (P2n11) at (6.642,1.641) {11};
\node[gn, fill=pc3, text=black] (P2n12) at (7.284,0.250) {12};
\node[ptitle] at (9.050,5.000) {(a) Stage 1 result};
\node[pstat] at (9.050,0.000) {scores on the edges of node 6: Jaccard\,|\,motif};
\node[barrow, rotate=0, minimum height=9mm] at (13.000,2.500) {};
\node[note, anchor=south] at (13.000,2.780) {motif};
\draw[goodc, line width=1.1pt, rounded corners=2pt] (13.650,0.000) rectangle (18.100,5.000);
\draw[ge] (14.939,4.259) -- (15.667,3.605);
\draw[ge] (14.939,4.259) -- (17.174,3.932);
\draw[ge] (14.939,4.259) -- (15.979,4.750);
\draw[ge] (14.939,4.259) -- (14.939,2.295);
\draw[ge] (15.667,3.605) -- (17.850,2.459);
\draw[ge] (15.667,3.605) -- (17.174,3.932);
\draw[ge] (15.667,3.605) -- (15.979,4.750);
\draw[ge] (15.667,3.605) -- (14.939,2.295);
\draw[ge] (17.850,2.459) -- (17.174,3.932);
\draw[ge] (17.850,2.459) -- (16.291,2.295);
\draw[ge] (17.174,3.932) -- (15.979,4.750);
\draw[ge] (17.174,3.932) -- (16.291,2.295);
\draw[ge] (17.174,3.932) -- (17.174,0.659);
\draw[ge] (16.291,2.295) -- (17.174,0.659);
\draw[ge] (16.291,2.295) -- (14.939,2.295);
\draw[ge] (17.174,0.659) -- (14.939,2.295);
\draw[ge] (14.939,2.295) -- (14.316,1.641);
\draw[ge] (14.108,0.659) -- (14.316,1.641);
\draw[ge] (14.108,0.659) -- (14.732,0.250);
\draw[ge] (13.900,2.459) -- (14.316,1.641);
\node[ring] at (16.291,2.295) {};
\node[gn, fill=pc0, text=white] (P3n1) at (14.939,4.259) {1};
\node[gn, fill=pc0, text=white] (P3n2) at (15.667,3.605) {2};
\node[gn, fill=pc0, text=white] (P3n3) at (17.850,2.459) {3};
\node[gn, fill=pc0, text=white] (P3n4) at (17.174,3.932) {4};
\node[gn, fill=pc0, text=white] (P3n5) at (15.979,4.750) {5};
\node[gn, fill=pc0, text=white] (P3n6) at (16.291,2.295) {6};
\node[gn, fill=pc0, text=white] (P3n7) at (17.174,0.659) {7};
\node[gn, fill=pc0, text=white] (P3n8) at (14.939,2.295) {8};
\node[gn, fill=pc1, text=white] (P3n9) at (14.108,0.659) {9};
\node[gn, fill=pc1, text=white] (P3n10) at (13.900,2.459) {10};
\node[gn, fill=pc1, text=white] (P3n11) at (14.316,1.641) {11};
\node[gn, fill=pc1, text=white] (P3n12) at (14.732,0.250) {12};
\node[ptitle] at (15.875,5.000) {(c) Stage 2 with motif similarity};
\node[pstat] at (15.875,0.000) {2 comm., NMI 1.000};
\end{tikzpicture}
\caption{Jaccard versus motif similarity, both run from the Stage~1 partition (a) of a planted graph (nodes 1--8 and 9--12).}
\label{fig:m1}
\end{figure*}
